\section{Computational study}
\label{sec:computational}

This section compares ALIC with the exact method of Prerna and
Sharma~\cite{prernaSharma2024}. The aim is to evaluate two complementary
aspects of the algorithms: the time required to certify an optimal efficient
solution and the number of optimization subproblems solved before
termination. Solution agreement is checked first, because both procedures
are exact.

The experiments were implemented in Python~3 in an Anaconda environment
and executed on a Windows workstation equipped with an 11th-generation
Intel Core i7 processor and 16~GB of RAM. All linear integer and
mixed-integer linear models were submitted through
\texttt{scipy.optimize.milp}, which uses the HiGHS optimization
backend~\cite{virtanen2020scipy,huangfuHall2018}. In ALIC, the exact
linearized master problem and the non-affine secondary completion are
MILPs. The primary completion/efficiency problem and the affine secondary
completion are linear integer models solved through the same interface.
For the Prerna--Sharma reconstruction, the same solver interface was used
for the preprocessing ILPs, the ranked linear problems, the tied-level
enumeration problems, and the efficiency tests. Using one solver backend
for both methods avoids a solver-dependent advantage.

The comparison follows a paired design. Each generated instance is solved
once by ALIC and once by Prerna--Sharma, and the execution order is
alternated from one instance to the next. The wall-clock time includes the
method-specific preprocessing. Internal LP relaxations performed by HiGHS
are not counted as separate subproblems. The source code, deterministic
instance generator and manifests, and archived result data are available at
\url{https://github.com/Ramdani-sudo/ALIC-Indefinite-Quadratic-Optimization-over-Integer-Efficient-Sets}.

\subsection{Benchmark generation and experimental protocol}
\label{subsec:benchmark-protocol}

The random generator follows the coefficient ranges used by Prerna and
Sharma~\cite{prernaSharma2024}. For every raw draw, the entries of the
constraint matrix are sampled from $\{0,\ldots,20\}$ and the right-hand side
from $\{0,\ldots,100\}$. The coefficients of the objective matrix and of
the quadratic preference are sampled from
$\{-100,\ldots,100\}$. The quadratic matrix is generated symmetrically and
is required to be strictly indefinite. The lower bounds are zero.

Because the constraint coefficients and the right-hand side are
nonnegative, finite integer upper bounds are obtained without solving
additional optimization problems:
\[
 u_j=\min_{i:A_{ij}>0}\left\lfloor\frac{b_i}{A_{ij}}\right\rfloor .
\]
This bound computation is part of instance construction and is identical
for both methods.

A direct use of the original random ranges produced many uninformative
draws. In particular, some instances had no nonzero unit feasible point,
and the quadratic preference could be constant on a simple probe set. We
therefore used a complete draw-and-reject procedure. A raw draw was kept
only when at least one nonzero unit vector was feasible, the preference was
nonconstant on the probe set formed by the origin and the feasible unit
vectors, and the quadratic matrix was strictly indefinite. The accepted
instance was never modified after generation. In total, 7614 raw draws
were rejected before the final benchmark set was obtained. These
restrictions remove obvious trivialities while preserving the coefficient
ranges of the Prerna--Sharma protocol.

The benchmark grid is
\[
 p\in\{5,10,20,50\},\qquad
 n\in\{10,20,50\},\qquad
 m\in\{10,20,30,50\},
\]
where $p$ denotes the number of criteria in the experimental generator.
There are $4\times3\times4=48$ configurations. Each configuration is
repeated 30 times, giving 1440 paired instances and 2880 method runs. A
time limit of 300~s is imposed on each method for each instance. A run is
declared optimal only when the solver certificate required by the
corresponding algorithm is available.

\subsection{Statistical methodology}
\label{subsec:statistics}

The observations are paired because both methods solve exactly the same
instances. Following modern recommendations for optimization
benchmarks~\cite{carrasco2020,bartzBeielstein2020}, the analysis combines
descriptive statistics, nonparametric inference, effect sizes, and a
Bayesian practical comparison.

For wall-clock time and solver calls, the main frequentist comparison is
the two-sided Wilcoxon signed-rank test~\cite{wilcoxon1945}. The
rank-biserial correlation is reported as an effect size. Because runtime
ratios are skewed, the Wilcoxon analysis is complemented by a paired
sign-permutation test on
$\log(T_{\mathrm{P}}/T_{\mathrm{A}})$~\cite{pesarin2001}. Bootstrap
95\% confidence intervals are used for the median speedup.

The 48 configuration-wise runtime tests form one family. Their $p$-values
are therefore adjusted with Holm's sequential procedure~\cite{holm1979}.
Optimal-solve success is compared with the exact McNemar test. In
addition, a Bayesian signed-rank comparison based on a Dirichlet process
is used~\cite{benavoli2014}. The practical-equivalence region is set to
$\pm10\%$ in relative runtime. Finally, a Dolan--Mor\'e performance
profile is used as a solver-oriented graphical summary~\cite{dolanMore2002}.

\subsection{Results and discussion}
\label{subsec:results-discussion}

The two methods reached certified optimality together on 1421 instances.
Their optimal objective values agree on all 1421 cases. The decision
vectors are identical in 1419 cases. In the remaining two cases, different
efficient decisions attain the same optimal quadratic value. Thus, no
instance produced two different certified optimal objective values.

ALIC solved 1438 of the 1440 instances, whereas the Prerna--Sharma
reconstruction solved 1421. ALIC was the only method to certify optimality
on 17 instances; the reverse case did not occur. The exact McNemar test
gives $p=1.53\times10^{-5}$.

\begin{table}[t]
\centering
\caption{Overall comparison on the complete benchmark. Runtime summaries
for the first four rows use the 1421 instances solved optimally by both
methods.}
\label{tab:overall-results}
\small
\begin{tabular}{lrr}
\toprule
Measure & ALIC & Prerna--Sharma\\
\midrule
Optimal solves (out of 1440) & 1438 & 1421\\
Median runtime (s) & 0.433 & 0.168\\
Mean runtime (s) & 1.314 & 2.969\\
Maximum common-optimal runtime (s) & 99.493 & 283.294\\
PAR-2 mean runtime (s) & 2.587 & 10.846\\
Median solver calls & 2 & 25\\
Mean solver calls & 2.118 & 44.975\\
\bottomrule
\end{tabular}
\end{table}

The main structural difference appears in the number of subproblems.
Across the 1421 common optimal solves, the median is 2 calls for ALIC and
25 for Prerna--Sharma. The median ratio
$N_{\mathrm{P}}/N_{\mathrm{A}}$ is 12.5. ALIC uses fewer optimization
subproblems on every common optimal pair. The Wilcoxon test gives
$p=1.49\times10^{-234}$ with rank-biserial correlation equal to 1.

The global paired runtime comparison favors Prerna--Sharma for the typical
common-optimal instance: the Wilcoxon test gives
$p=4.39\times10^{-78}$ and a rank-biserial effect of $-0.573$ when a
positive sign is defined in favor of ALIC. The median speedup
$T_{\mathrm{P}}/T_{\mathrm{A}}$ is 0.237, with a 95\% bootstrap interval
of approximately $[0.221,0.256]$.

The mean runtime tells a different part of the story. It is 1.314~s for
ALIC and 2.969~s for Prerna--Sharma. The PAR-2 mean is also much smaller
for ALIC. These measures include the difficult tail and the timeout
penalty, and therefore show that ALIC is more robust on the hardest
instances even though it is not the fastest method on the median
common-optimal run.

The configuration-level analysis is also informative. After Holm
correction of the 48 paired Wilcoxon tests, 31 configurations
significantly favor Prerna--Sharma in median runtime, 4 significantly
favor ALIC, and 13 show no significant difference. The four significant
ALIC wins all occur for $n=50$ and $m=50$, for
$p\in\{5,10,20,50\}$. Their median speedups are between 6.52 and 7.03.
This crossover is consistent with the purpose of ALIC: its larger master
models can be costly on easy instances, but the strong reduction in the
number of repeated subproblems becomes valuable on difficult instances.

\subsection{Limitations and interpretation of the computational cost}

The experiments show that fewer optimization subproblems do not
necessarily imply a lower wall-clock time. ALIC solves substantially
fewer subproblems than the Prerna--Sharma method, but its main subproblem
is an exact master MILP containing the binary expansion variables,
auxiliary product variables, and the accumulated image-cut variables.
This master can therefore be more expensive than the smaller ILPs solved
by the reference method.

The primary completion of ALIC is lighter, since it is an ILP used to
test efficiency and to obtain an efficient dominating solution when
needed. A secondary problem is solved only when uniqueness is not
certified, and it is either an ILP or an exact MILP depending on the
restricted quadratic structure. In contrast, the Prerna--Sharma method
solves a larger sequence of ranking, tied-level enumeration, and
efficiency-test problems. Hence, its number of solver calls is much
larger, but many of these calls are small.

This explains the observed results. On the 1421 common optimal instances,
the median number of solver calls is 2 for ALIC and 25 for
Prerna--Sharma, while the median runtime is 0.433~s and 0.168~s,
respectively. Thus, Prerna--Sharma is often faster on easy instances.
However, ALIC performs better on the difficult tail: its mean runtime is
1.314~s versus 2.969~s, and its PAR-2 score is 2.587~s versus 10.846~s.
The main practical advantage of ALIC is therefore a much smaller number of
optimization subproblems and better robustness on difficult instances,
rather than a uniform runtime advantage.

Finally, the benchmark is synthetic and follows the Prerna--Sharma
generation ranges, with additional rejection rules used to remove obvious
trivial cases. The conclusions should therefore be restricted to the
tested setting.
